\documentclass[11pt,a4paper]{article}
\usepackage[utf8]{inputenc}
\usepackage[T1]{fontenc}
\usepackage[english]{babel}
\usepackage{amsmath,amssymb,amsthm}
\usepackage{graphicx}
\usepackage{hyperref}
\usepackage{geometry}
\geometry{margin=2.5cm}

\title{Theory of Universal Electro-Intelligence (TEIU)\\[2mm]
\large $I = \Phi \times \log(1+C) \times \frac{1}{1+D}$}
\author{Gregory Lombardo\\Independent Researcher}
\date{Version 1.1 --- 14 November 2026}

\begin{document}
\maketitle

\begin{abstract}
We propose the Theory of Universal Electro-Intelligence (TEIU): electricity is not only a carrier of energy, but the physical substrate in which intelligence can emerge, persist, and be measured.
We define electro-intelligence $I$ as a function of three measurable quantities: integration $\Phi$, complexity $C$, and dissipation $D$.
We propose a falsifiable formula, five testable predictions, and a first numerical test (Simulation 01).
Version 1.1 corrects Simulation 01: Gaussian $\Phi$ proxy, chaotic task, spectral-radius normalisation. $I$ now increases with $C$ and falls with $D$, as predicted.
Code and data: GitHub \url{https://github.com/tra3600/teiu}.
Archiving: Zenodo concept DOI \url{https://doi.org/10.5281/zenodo.24202580}, Version v1.1 DOI \url{https://doi.org/10.5281/zenodo.24203665}, Version v1.0.1 DOI \url{https://doi.org/10.5281/zenodo.24202581}.
\end{abstract}

\section{Introduction}
Electricity is everywhere: in neurons, in circuits, in plasmas, in the interstellar medium. Intelligence, as observed in brains and in artificial systems, always appears in electrically active matter. TEIU asks a simple question: is this a coincidence, or is electricity the universal substrate of intelligence?

We do not claim that all electricity is intelligent. We claim that intelligence, when it appears, appears as a property of electrically active, integrated, complex, low-dissipation systems, and that this property can be measured.

\section{Axioms}
\textbf{Axiom 1 (Substrate).} All observed intelligence emerges in electrically active matter. No counter-example of intelligence without electrical activity is known.

\textbf{Axiom 2 (Measurability).} Intelligence can be quantified by integration, complexity, and dissipation, all measurable in principle from time series and network structure.

\textbf{Axiom 3 (Falsifiability).} If a highly integrated, highly complex, low-dissipation electrical system systematically shows no predictive or adaptive capacity, while a poorly integrated system shows more, TEIU is falsified.

\section{The Formula}
We define electro-intelligence as:
\begin{equation}
I = \Phi \times \log(1+C) \times \frac{1}{1+D}
\label{eq:I}
\end{equation}
where:
\begin{itemize}
\item $\Phi$ is integration, the information generated by the whole beyond the sum of its parts, in the sense of Integrated Information Theory, approximated here by a Gaussian entropy proxy.
\item $C$ is complexity, combining mean degree, density, and clustering (triangles) of the network.
\item $D$ is dissipation, energy lost as heat / noise per unit of computation, here a control parameter in $[0,1]$.
\end{itemize}
$\log(1+C)$ prevents unbounded growth with size. $1/(1+D)$ penalises dissipation: high dissipation destroys intelligence even if $C$ is high.

\section{Predictions}
\textbf{P1.} $I$ increases with $C$ when $D$ is low and $\Phi>0$.

\textbf{P2.} $I$ collapses when $D$ is high, even if $C$ is high.

\textbf{P3.} A system with $\Phi \approx 0$ has $I \approx 0$, whatever its $C$.

\textbf{P4.} $I$ predicts adaptive capacity (prediction score $R^2$ on a non-trivial task) better than $C$ alone.

\textbf{P5.} In biological data, cortical networks should show higher $I$ than subcortical or random controls with matched $C$.

\section{Simulation 01 v1.1 --- Corrected}
\subsection{What was wrong in v1.0}
v1.0 used pure sinusoids (too easy, $R^2 \approx 0.99$ everywhere) and a binarised $\Phi$ proxy that collapsed to $0$ for large networks. $I$ was $0$ for 3 networks out of 4. The simulation did not test the theory.

\subsection{Corrections in v1.1}
\begin{itemize}
\item \textbf{Chaotic task:} logistic map $x_{t+1}=3.99\,x_t(1-x_t)$ + sinus + noise. No longer trivially predictable.
\item \textbf{Gaussian $\Phi$:} $\Phi = H(\text{left}) + H(\text{right}) - H(\text{total})$, with Gaussian entropy $H = \frac{1}{2}\log\det(2\pi e \Sigma)$. No binarisation. $\Phi$ no longer collapses to $0$ when $n$ grows.
\item \textbf{Spectral-radius normalisation:} $W$ rescaled to target spectral radius $\rho = 0.30$ (poor), $0.70$ (medium), $0.95$ (rich), with signed weights and leaky integration $X \leftarrow 0.7X + 0.3\tanh(\text{rec}+\text{ext}+\text{noise})$. Without this, the rich network saturates in $\tanh$ and loses memory.
\end{itemize}

\subsection{Results v1.1 --- Validated locally, identical on two machines}
\begin{center}
\begin{tabular}{lccccc}
\hline
Network & $\Phi$ & $C$ & $D$ & $R^2$ & $I$ \\
\hline
Poor & 3.082 & 6.05 & 0.35 & 0.253 & 4.458 \\
Medium & 3.325 & 28.21 & 0.15 & 0.246 & 9.755 \\
Rich & 3.273 & 73.31 & 0.08 & 0.244 & 13.056 \\
Rich, dissipative & 3.039 & 74.09 & 0.40 & 0.256 & 9.375 \\
\hline
\end{tabular}
\end{center}

Interpretation: $\Phi$ is now stable around $3$, no longer $0$. $R^2$ no longer saturates at $0.99$, it is around $0.25$, realistic for a chaotic target. $I$ follows the theory: $4.458 < 9.755 < 13.056$, and the rich dissipative network falls back to $9.375$. $D$ breaks $I$ even when $C$ is high. This is P1 and P2.

\subsection{Known limitation, published in v1.1}
$R^2$ remains flat around $0.25$ for all networks. The chaotic target is too hard for this small reservoir ($n=60$) to distinguish networks by $R^2$ alone. $I$ distinguishes them via $C$ and $D$, but $R^2$ alone does not. v1.2 will use an intermediate task (neither too easy nor too hard) so that $R^2$ also distinguishes networks.

Reproduction:
\begin{verbatim}
pip install -r requirements.txt
python simulations/teiu_simulation.py --plot
\end{verbatim}
Files: \texttt{simulations/teiu\_simulation.py}, \texttt{figures/teiu\_sim01\_v1.1\_result.png}.

\section{Magnetism and Self-Consciousness --- v2.0 Draft, Not in v1.1}
For v2.0, we will test $B = \nabla \times W$, a magnetic-like field derived from the weight field $W$, and a coherence measure $B_{\text{coh}}$. Hypothesis: intelligence without self-consciousness is $I$ as in Eq.~\ref{eq:I}. Intelligence with self-consciousness is $I_{\text{with-self}} = I \times (1 + B_{\text{coh}})$. This is a draft, not claimed in v1.1, to be tested in Simulation 04 before integration.

\section{Conclusion}
TEIU v1.1 is a corrected, falsifiable first step. It does not prove that electricity contains all intelligence. It provides a measurable $I$, a testable prediction, and a simulation that now behaves as predicted, with its limits published. The next step is independent replication, biological tests (P5), and v1.2 with an intermediate task.

\section*{Data Availability}
GitHub: \url{https://github.com/tra3600/teiu} --- Release v1.1: \url{https://github.com/tra3600/teiu/releases/tag/v1.1}\\
Zenodo concept DOI (all versions): \url{https://doi.org/10.5281/zenodo.24202580}\\
Zenodo v1.1: \url{https://doi.org/10.5281/zenodo.24203665}\\
Zenodo v1.0.1: \url{https://doi.org/10.5281/zenodo.24202581}\\
Licences: Code MIT, Texts CC BY 4.0.

\end{document}
